% Generated draft: requires mathematical and source review.
\documentclass[11pt]{article}
\usepackage[margin=25mm]{geometry}
\usepackage{fontspec}
\setmainfont{Latin Modern Roman}
\usepackage{amsmath,amssymb,mathtools,unicode-math}
\setmathfont{Latin Modern Math}
\usepackage{enumitem,needspace,xurl,hyperref}
\setlength{\parindent}{0pt}
\setlength{\parskip}{5pt}
\setlength{\emergencystretch}{3em}
\newcommand{\N}{\mathbb{N}}
\newcommand{\Z}{\mathbb{Z}}
\newcommand{\Q}{\mathbb{Q}}
\newcommand{\R}{\mathbb{R}}
\newcommand{\C}{\mathbb{C}}
\newcommand{\F}{\mathbb{F}}
\newcommand{\A}{\mathbb{A}}
\newcommand{\PP}{\mathbb{P}}
\newcommand{\E}{\mathbb{E}}
\newcommand{\eps}{\varepsilon}
\newcommand{\dd}{\,\mathrm{d}}
\begin{document}
\section*{UnsolvedMath ID: 30005681}
\textbf{Source problem code:} catalogue entry
\subsection{Definitions and mathematical statement}
A group \(G\) is \emph{residually finite} if, for every \(g\in G\setminus\{e\}\), there exist a finite group \(F\) and a homomorphism \(\varphi:G\to F\) such that \(\varphi(g)\ne e_F\). Equivalently,
\[
\bigcap_{\substack{N\trianglelefteq G\\ [G:N]<\infty}}N=\{e\}.
\]
The supplied record uses the term \emph{chaotic almost minimal action}, abbreviated CAM, without defining it. Write \(\mathcal{C}(G)\) provisionally for the assertion that \(G\) admits an action of the type called CAM in the cited report; this notation is a placeholder for the missing definition.

The recorded existence question is:
\begin{quote}
Does every residually finite group \(G\) admit a chaotic almost minimal action?
\end{quote}
Its logical form is
\[
\forall G\;\bigl(G\text{ is a residually finite group}\Longrightarrow\mathcal{C}(G)\bigr).
\]
The supplied report passage also asserts the necessary implication
\[
\mathcal{C}(G)\Longrightarrow G\text{ is residually finite}.
\]
This implication is reported source context, not a theorem verified in this draft.

\textbf{Unresolved qualifications.} The record does not supply the definition of CAM, the class of underlying spaces, or the requirements on the action. It also does not state whether the groups must be infinite, countable, or finitely generated. The literal universal wording includes finite groups and the trivial group; whether these are intended must be checked against the defining source. Thus the displayed question preserves the supplied quantifier, but a fully specified mathematical formulation requires the missing definition and scope restrictions.
\subsection{Short English statement}
The question asks whether residual finiteness is sufficient for a group to admit a chaotic almost minimal action. The supplied passage states that residual finiteness is necessary. A precise interpretation still requires the definition of CAM and clarification of the permitted groups and action spaces.
\subsection{Sources}
\begin{itemize}
\item Supplied local UnsolvedMath registry record, permanent ID \(30005681\), entitled \emph{CAM Actions of Residually Finite Groups}; exact source problem code \texttt{catalogue entry}. The record lists proposed year 2023 and supplies no named proposer.
\item Supplied source citation: \emph{Aspects of Aperiodic Order (2024)}, Oberwolfach Reports, DOI \texttt{10.4171/owr/2023/37}, pp.\ 2107--2178. \url{https://doi.org/10.4171/owr/2023/37}. The background describes Oberwolfach Reports as documenting mathematical workshops and their open-problem discussions.
\item Supplied literature source labelled ``Oberwolfach report'': \url{https://doi.org/10.4171/owr/2023/37}.
\item Supplied literature source labelled ``Cyr-Kra-Schmieding'': \url{https://sites.math.northwestern.edu/~kra/papers/CAM.pdf}. No further bibliographic details are supplied.
\item Literature assessment dated \texttt{2026-08-22}: the supplied curation labels the problem open, states that the report presents the existence question as open, and reports that no later primary result located by its cutoff resolves it. This is supplied dated curation, not an independently verified literature assessment.
\item The registry labels the statement \texttt{corrected\_verified} and says it was reconstructed and checked against the official report. These are supplied verification claims; the report and linked paper have not been independently consulted for this draft.
\end{itemize}
\end{document}
